\chapter{Tutorial on setting the Research Details}\label{chap:param}

The preamble of your research book will contain the cover page, abstract of the research and table of contents. In these parts, several things (the title of the research, name of the students and the supervisor etc.) will appear multiple times. So, we have set up an arrangement so that you can write these details or parameters once and have them written in different places at once.

\section{Information Folder}
\textbf{Information} folder in the template contains several .txt files consisting of the research details.
\begin{description}

\item [date.txt] This file contains the month and year of the research book. Change it and compile the project to see the changes in cover page.

\item [students.txt] This file contains the name and IDs of the students. There are one name per line. Follow the name with the student ID. Name and ID must be separated by a comma. Change the names and IDs and compile the template to see the changes accordingly in cover page.

\item [supervisor.txt] This file contains the name and the designation of the Supervisor of the research group. Name and designation will be separated by a comma. Change them and compile the project to see the changes in cover page.

\item [title.txt] This file contains the name of the research. Change it and compile the project to see the changes in cover page. 

\item [worktype.txt] This file contains the type of your work (research). Ignore it fr now.
\end{description}

\section{Prologue Folder}
\textbf{Prologue} folder contains large texts to be changed- the likes of abstract. There is \textbf{abstract.tex} file where you can write the corresponding texts in \LaTeX and see the changes in abstract page. 
\endinput
